% ============================================================================
% Fichier  : partie4/P04_C19_reprise.tex
% Titre    : Reprise Informatique et Résilience
% Auteur   : MINKA MI NGUIDJOÏ Thierry Emmanuel
% Version  : 1.0 -- novembre 2026
% ============================================================================
% Sources : ISO/IEC 27031:2011 ; NIST SP 800-34 Rev. 1 (2010) ; Comité de
% Bâle, Principles for Operational Resilience (2021). Chiffres de LiZbiZ :
% cas fictif, cohérents avec SEC3266 (arrêt de 250 h 20 de la facturation).
% ============================================================================

\chapter{Reprise Informatique et Résilience}
\label{chap:pra}
\chaptermark{Reprise informatique et résilience}

\epigraph{%
    « Après un rançongiciel, on ne restaure pas un système~;\\
    on en reconstruit un sain. »%
}{Principe du cours}

\niveauChapitre{Intermédiaire}{Chapitres~\ref{chap:bia}
et~\ref{chap:pca}~; cours SEC3266, chapitre sur la continuité.}

\begin{boxobjectifs}
\begin{itemize}
    \item Choisir une technique de protection des données selon le RTO et
          le RPO.
    \item Rédiger un plan de reprise informatique.
    \item Conduire une reprise après rançongiciel sans réinfecter.
    \item Appliquer les principes de résilience opérationnelle.
    \item Mesurer l'écart entre objectifs et capacité réelle.
\end{itemize}
\end{boxobjectifs}

\bigskip

% ============================================================================
\section*{Question directrice : comment passer de 250 heures à moins de 48~?}
% ============================================================================

En mars, la facturation de LiZbiZ est restée arrêtée 250~h~20, pour un
objectif fixé après coup à 48~heures. En juin, M.~KOUNOU doit présenter un
plan de reprise qui tienne cet objectif, et le prouver par un test avant
la fin de l'année.

\medskip
\textbf{Plan du chapitre.}
\begin{enumerate}
    \item Protéger les données selon les objectifs (section~\ref{sec:c19-techniques}).
    \item Le plan de reprise informatique (section~\ref{sec:c19-pra}).
    \item Reprendre après un rançongiciel (section~\ref{sec:c19-rancon}).
    \item La résilience opérationnelle (section~\ref{sec:c19-resilience}).
    \item Mesurer la capacité réelle (section~\ref{sec:c19-mesurer}).
\end{enumerate}

\separateur

% ============================================================================
\section{Protéger les Données selon les Objectifs}
\label{sec:c19-techniques}
% ============================================================================

\begin{tableaucours}{L{3.6cm}L{2.4cm}L{2.4cm}L{3.8cm}}{%
    \tableauHeader{Technique} &
    \tableauHeader{RPO typique} &
    \tableauHeader{RTO typique} &
    \tableauHeader{Limite}
}
Sauvegarde et restauration & Heures à 1~jour & Heures à jours & Lente à
restaurer en volume \\
Réplication asynchrone & Minutes & Heures & Réplique aussi le chiffrement
malveillant \\
Réplication synchrone & Nul & Minutes à heures & Distance limitée, coût \\
Actif-actif sur deux sites & Nul & Minutes & Coût élevé, complexité \\
\end{tableaucours}
\vspace{-0.8em}
\begin{center}{\small\itshape Tableau~19.1. Techniques de protection des données}\end{center}

\begin{boxalert}
\textbf{Violation typique.} Compter sur la réplication contre le
rançongiciel.

\smallskip
\textbf{Règle~:} la réplication protège contre la perte d'un site, pas
contre la corruption~: elle recopie le chiffrement. Contre le
rançongiciel, il faut des sauvegardes \emph{versionnées}, dont une copie
\emph{hors ligne ou immuable}, administrée hors du domaine, et testée.
\end{boxalert}

% ============================================================================
\section{Le Plan de Reprise Informatique}
\label{sec:c19-pra}
% ============================================================================

L'\norme{ISO/IEC~27031} traite de la préparation des technologies de
l'information à la continuité~\cite{iso27031_2011}~; le NIST SP~800-34 décrit
les plans de contingence des systèmes~\cite{nist80034}.

\begin{boxprocedure}
\textbf{Contenu du plan de reprise informatique de LiZbiZ}
\begin{enumerate}
    \item Inventaire des systèmes, avec leur RTO et leur RPO (issus de
          l'analyse d'impact).
    \item Ordre de reprise selon les dépendances~: réseau, annuaire, bases,
          applications.
    \item Procédures de restauration, pas à pas, pour chaque système.
    \item Site de reprise (centre de Yaoundé), capacités, accès.
    \item Rôles~: équipe informatique, prestataire, éditeurs.
    \item Critères de réussite et vérifications (intégrité des données,
          tests fonctionnels avec les métiers).
    \item Retour sur le site principal.
\end{enumerate}
\end{boxprocedure}

% ============================================================================
\section{Reprendre après un Rançongiciel}
\label{sec:c19-rancon}
% ============================================================================

Après un rançongiciel, l'attaquant peut avoir laissé des accès (comptes
créés, services installés). Restaurer sur l'infrastructure compromise, ou
restaurer une sauvegarde postérieure à son entrée, revient à le réinviter.

\begin{boxprocedure}
\textbf{Reprise après rançongiciel}
\begin{enumerate}
    \item Préserver les preuves avant toute remise en état
          (cours SEC3266).
    \item Établir la date d'entrée de l'attaquant~; choisir un point de
          restauration antérieur, ou vérifier la sauvegarde.
    \item Reconstruire un environnement sain, isolé~: annuaire neuf, comptes
          réinitialisés depuis des postes sains.
    \item Restaurer les données, pas les systèmes infectés~; réinstaller les
          systèmes à partir de sources sûres.
    \item Vérifier l'absence de persistance (comptes, services, tâches
          planifiées) avant de reconnecter.
    \item Reconnecter progressivement, sous surveillance renforcée.
\end{enumerate}
\end{boxprocedure}

\begin{boxscenario}{Opération GAROUA. Pourquoi onze jours}
\textbf{Constatations.} Sauvegardes en ligne supprimées le 5~mars à
21~h~15~; restauration à partir d'une copie de fin janvier retrouvée sur un
disque~; ressaisie de six semaines de factures~; annuaire reconstruit du
9 au 12~mars.

\medskip
\textbf{Analyse} (confiance élevée). Trois causes expliquent l'essentiel de
l'arrêt~: l'absence de copie immuable (six semaines de données à
ressaisir), l'absence de procédure de reconstruction de l'annuaire (quatre
jours), et l'absence de site de reprise prêt. Le plan de juin traite les
trois.
\end{boxscenario}

% ============================================================================
\section{La Résilience Opérationnelle}
\label{sec:c19-resilience}
% ============================================================================

\begin{boxdef}
\termecle{Résilience opérationnelle}~: capacité d'une organisation à
fournir ses opérations critiques à travers les perturbations, en anticipant,
en absorbant, en s'adaptant et en se rétablissant.
\end{boxdef}

Le Comité de Bâle sur le contrôle bancaire a publié en 2021 des principes
de résilience opérationnelle~\cite{bcbs2021resilience}. Ils déplacent le
regard des systèmes vers les \emph{opérations critiques}~: pour chacune, on
fixe une \emph{tolérance à la perturbation}, on cartographie ses
interconnexions (personnes, technologies, tiers), et l'on teste la
capacité à rester dans la tolérance face à des scénarios sévères mais
plausibles.

\begin{tableaucours}{L{4.8cm}L{7.4cm}}{%
    \tableauHeader{Principe} &
    \tableauHeader{Application à la BANQUE SANAGA}
}
Gouvernance & Le conseil approuve les tolérances \\
Cartographie & Virements~: système central, annuaire, connexion aux
paiements, équipes de Douala et de Yaoundé \\
Gestion des tiers & Prestataire monétique, fournisseur du système central \\
Continuité testée & Bascule annuelle sur le site de secours \\
Gestion des incidents & Grille et escalade adaptées aux virements \\
\end{tableaucours}
\vspace{-0.8em}
\begin{center}{\small\itshape Tableau~19.2. Principes de résilience opérationnelle (d'après le Comité de Bâle, 2021)}\end{center}

% ============================================================================
\section{Mesurer la Capacité Réelle}
\label{sec:c19-mesurer}
% ============================================================================

\begin{boxobjet}{R-01 · Test de reprise de novembre}
\textbf{Scénario testé~:} perte du centre de Douala~; reprise de la
facturation à Yaoundé à partir de la sauvegarde immuable de la nuit.
\textbf{Résultat~:} facturation rétablie en 19~heures (RTO de 48~heures
tenu)~; données perdues~: 5~heures (RPO de 4~heures non tenu, la
sauvegarde ayant eu lieu à 1~h et la panne simulée à 6~h).
\textbf{Action~:} journalisation des transactions entre deux sauvegardes.
\textbf{Propriétaire~:} M.~KOUNOU. \textbf{Indicateur~:} durée de la dernière
reprise mesurée.
\end{boxobjet}

\begin{boxmetier}{Responsable de la continuité}
\textbf{Comparer la capacité mesurée aux objectifs.} Le responsable tient,
pour chaque système critique, deux colonnes~: l'objectif (RTO, RPO) et la
dernière valeur mesurée en test. Un objectif jamais mesuré est noté
« inconnu », non « atteint ».
\end{boxmetier}

\separateur

\begin{boxresume}
\textbf{Ce qu'il faut retenir de ce chapitre~:}
\begin{itemize}
    \item \textbf{La technique} se choisit selon le RTO et le RPO~; la
          réplication ne protège pas de la corruption.
    \item \textbf{Le plan de reprise} fixe l'ordre, les procédures, le
          site, les rôles et les critères de réussite.
    \item \textbf{Après un rançongiciel}, on reconstruit un environnement
          sain et l'on restaure les données, pas les systèmes infectés.
    \item \textbf{La résilience opérationnelle} raisonne en opérations
          critiques et en tolérances.
    \item \textbf{La capacité réelle} se mesure en test et se compare aux
          objectifs.
\end{itemize}
\end{boxresume}

\begin{boxinfo}
\textbf{Points de contrôle de l'élève.} Je sais~: associer une technique à
des objectifs~; structurer un plan de reprise~; dérouler une reprise après
rançongiciel~; énoncer les principes de résilience opérationnelle~;
interpréter un résultat de test.
\end{boxinfo}

\bigskip

\begin{boxexo}[title={\ding{45}\quad Exercice~19.1. Quelle technique~? \niveaubase}]
Quelle technique du Tableau~19.1 choisir pour les virements de la BANQUE
SANAGA (RPO nul, RTO 2~h)~? Et pour les rapports réglementaires~?
\end{boxexo}

\begin{boxexo}[title={\ding{45}\quad Exercice~19.2. Le test \niveaubase}]
Dans le test de novembre, le RTO et le RPO sont-ils tenus~? Justifiez par
les chiffres.
\end{boxexo}

\begin{boxexo}[title={\ding{45}\quad Exercice~19.3. Réplication \niveaubase}]
Pourquoi la réplication n'aurait-elle pas sauvé LiZbiZ en mars~?
\end{boxexo}

\begin{boxexo}[title={\ding{45}\quad Exercice~19.4. Le point de restauration \niveaumoyen}]
L'attaquant est entré le 2~mars à 23~h~12. Quelle sauvegarde faut-il
restaurer, et à quelles vérifications la soumettre si l'on en prend une
plus récente~?
\end{boxexo}

\begin{boxexo}[title={\ding{45}\quad Exercice~19.5. L'ordre \niveaumoyen}]
Établissez l'ordre de reprise des systèmes de LiZbiZ~: réseau, annuaire,
base de facturation, application de facturation, messagerie, portail.
\end{boxexo}

\begin{boxexo}[title={\ding{45}\quad Exercice~19.6. Persistance \niveaumoyen}]
Citez quatre traces de persistance à rechercher avant de reconnecter un
système restauré.
\end{boxexo}

\begin{boxexo}[title={\ding{45}\quad Exercice~19.7. Tolérance \niveaumoyen}]
Formulez la tolérance à la perturbation des virements de la BANQUE
SANAGA, au sens du Comité de Bâle.
\end{boxexo}

\begin{boxexo}[title={\ding{45}\quad Exercice~19.8. Inconnu \niveaumoyen}]
Pourquoi noter « inconnu » un objectif jamais mesuré~? Reliez à l'instrument
du Chapitre~\ref{chap:retex}.
\end{boxexo}

\begin{boxexo}[title={\ding{45}\quad Exercice~19.9. Un plan de reprise \niveauexpert}]
Rédigez la procédure de reprise de la facturation de LiZbiZ sur le site de
Yaoundé, pas à pas, avec les vérifications.
\end{boxexo}

\begin{boxexo}[title={\ding{45}\quad Exercice~19.10. Scénario sévère \niveauexpert}]
Proposez à la BANQUE SANAGA trois scénarios « sévères mais plausibles »
pour tester sa résilience, et la mesure de réussite de chacun.
\end{boxexo}

\bigskip

\begin{boxinfo}
\textbf{Pour aller plus loin.}
\begin{itemize}
    \item ISO/IEC~27031~\cite{iso27031_2011}~; NIST SP~800-34~\cite{nist80034}.
    \item Comité de Bâle, \emph{Principles for Operational
          Resilience}~\cite{bcbs2021resilience}.
\end{itemize}
\end{boxinfo}

\begin{boxinfo}
\textbf{Transition.} Un plan non exercé est une hypothèse. Le
Chapitre~\ref{chap:exercer} apprend à exercer et à tester les plans.
\end{boxinfo}
